\section{Recoil-point-cloud representation cross-check}
\label{sec:pet}

\bpn{left off here}

\paragraph{Status: diagnostic and method development, not a publication
uncertainty product.}
Everything in this section is reported as a diagnostic and method-development
result. No PET covariance --- historical recoil-only or full-event --- is adopted
as a publication uncertainty product, and the full-event PET
central/statistical pairing was declined on 2026-08-20 rather than resolved: the
corresponding statistical covariance remains constructed but not independently
verified, and it is not paired with a PET central value. That decision does not
claim a PET central value is wrong and does not invalidate the dispersion
measured from a replica family; the missing evidence is coverage. No mechanism
for the anomaly it rests on is established, and none is asserted here.
Reconsideration requires estimator equivalence \emph{and} coverage evidence, and
coverage is a different object from checking how a replica covariance was
assembled --- independently verifying the construction would not satisfy that
gate. The publication uncertainty critical path runs through the adopted scalar
five-dimensional covariance (\S\ref{sec:uq-combine}), not through this section.
Routed record: \texttt{OI-126} in \texttt{docs/OPEN\_ITEMS.md}.

As a low-level-representation cross-check, a point-cloud transformer (PET, the architecture of
OmniLearn~\cite{Mikuni:2024qsr}, vendored with the upstream OmniFold
implementation) unfolds the same ME~FHC data using, on the reconstructed
side, the actual calorimetric recoil clusters --- per-cluster energy,
one view-projected transverse coordinate and $z$ (\texttt{cluster\_energy},
\texttt{cluster\_pos}, \texttt{cluster\_z}, filtered on
\texttt{cluster\_isMuontrack} via
\texttt{CVUniverse::GetRecoClusters()}). At truth level it receives the
final-state hadron cloud with the muon removed. Muon-derived event scalars are
retained for selection and binning the extracted result, but they are not
classifier features at either level. The learned weight is therefore a
function of the recoil representation, not of the complete joint muon--hadron
event. We therefore compare its absolutely normalized central value, closure
and replica dispersion only as diagnostic properties of the recoil-only
estimator. A full-event
estimator would instead train jointly on the declared muon and recoil inputs.

\subsection{From shape check to absolute cross section}
\label{sec:pet-abs}
The PET network trains on a GPU on a 2M-event subsample (2 OmniFold
iterations); the learned step-2 truth reweight is then applied to the
\emph{full} 32.8M-event generator sample
(\texttt{MultiFold.reweight}), binned on the 4D grid, and pushed through the
standard extraction (\texttt{extract\_cross\_section\_nd} with the GBDT
run's completeness --- completeness is reweight-independent --- and the same
flux/POT/nucleon constants). Three gates validate the absolute path:
\begin{enumerate}[itemsep=2pt]
  \item \textbf{Full-statistics reweight sanity}: the pushed weights over the
        32.8M generator events are finite with mean $1.028$;
  \item \textbf{Ordinary closure}: unfolding MC-reco-as-pseudo-data
        recovers MC truth to total $0.9884$ with uniform per-axis median
        residuals of $\sim\SI{1.1}{\percent}$ --- the absolute machinery is
        internally consistent. Because pseudo-data and prior come from the
        same simulation, this test cannot expose omitted muon dependence at
        fixed recoil;
  \item \textbf{Cross-check vs GBDT} (\textbf{\texttt{niter}${}={}$2 legacy --- not
        a current number}; see the scope note below): on data, the PET absolute total was
        $\dead{\SI{2.796e-38}{cm^2/nucleon}}$ versus the GBDT
        $\dead{\SI{3.066e-38}{}}$, ratio $\dead{\petRatio}$, with per-axis median
        differences of \dead{\SIrange{6.5}{9.9}{\percent}} (Fig.~\ref{fig:petabs}).
\end{enumerate}

\paragraph{Scope of the GBDT cross-check: \texttt{niter}${}={}$2 legacy, withdrawn
as a current comparison.}
Every struck figure in item~3 above, and every quantity derived from them, is
retained here only as a record of what the two-iteration run gave. They are
\textbf{not} the current PET/GBDT comparison, for reasons that are mechanical
rather than matters of judgement:
\begin{itemize}[nosep]
  \item \textbf{Iteration count.} The PET run used a 2M-event, \emph{two}-iteration
        training. The campaign's pinned policy is \texttt{niter}${}={}$3 (CLM-010,
        CLM-012, \textsc{frozen}), and full-event PET moved to it. The comparison
        has not been re-extracted at \texttt{niter}${}={}$3.
  \item \textbf{Estimator definition.} The full-event schema landed 2026-08-01 and
        the loader now reads the whole event, so a pre-08-01 PET total is a
        different estimator, not a differently-trained one.
  \item \textbf{Background subtraction.} The quoted 4D ratio was taken with unit
        measured weights and \emph{no} background subtraction (audit finding J21).
        At this analysis's ${\sim}\SI{3}{\percent}$ background scale, leaving
        background in biases PET high, so the underlying gap is \emph{larger} than
        the struck value, by an amount that is a first-order estimate and
        deliberately not quoted.
\end{itemize}
Note that the struck ratio is \emph{arithmetically correct} --- it is exactly the
quotient of the two struck operands. It is withdrawn because its operands describe
a training configuration that no longer exists, which is a failure a consistency
check cannot detect. Both operands are inline literals rather than managed macros,
so they are marked here at the point of use; the macro itself is flagged in
\texttt{values.tex}. See
\texttt{docs/orchestration/INDEX-retracted-and-superseded-values.md}.

What survives independently of the withdrawn comparison: the
$\dead{\sim\SI{\petGbdtGap}{\percent}}$ data-side gap is not
a normalization bug, but ordinary closure alone cannot assign it uniquely to
training configuration. A higher-epoch retrain moved the PET total by
$\dead{-1.6\%}$ (ratio $\dead{0.897}$) --- struck on the same grounds, being the
same legacy run --- which demonstrated training sensitivity qualitatively. The
recoil-only feature contract is an additional structural difference: at fixed
recoil, the conditional muon distribution remains inherited from the prior.

\begin{figure}[H]
  \centering
  \includegraphics[width=0.95\textwidth]{pet_vs_gbdt_absolute}
  \caption{Recoil-only PET representation cross-check: absolutely-normalized
  4D cross section from the point-cloud unfold (full-statistics reweight of the generator sample) versus the
  production scalar GBDT result, projected on each axis. The PET/GBDT total
  ratio is $\dead{\petRatio}$ --- \textbf{struck: \texttt{niter}${}={}$2 legacy, see
  the scope note in the text}; ordinary closure validates the extraction machinery to
  $\sim\SI{\petClosure}{\percent}$ but does not test omitted muon dependence.}
  \label{fig:petabs}
\end{figure}

\subsection{Statistical and training diagnostics of the recoil-only estimator}
\label{sec:pet-syst}
The PET statistical replicas independently
Poisson-fluctuate data and MC, retrain PET at a fixed estimator seed, and apply
the same coherent MC draw in the final truth binning and completeness.

The calculation uses the same background-subtracted five-dimensional
measured target as the scalar GBDT. Event-by-event data alignment is exact, the
MC arrays are byte-identical to the full-cloud input. The
nominal PET extraction, an identical-seed GPU-floor repeat, a
20-member coherent-Poisson retraining ensemble, and a 12-member crossed
subsample/estimator-seed ensemble have all passed their ordering, event-alignment, and
finite-output gates. The GPU nondeterminism is negligible relative to both
sampling and training variation. The $C_{\rm stat}$ block is a diagnostic
dispersion measurement built from 20 coherent replicas. It is not an
independently verified covariance, it is not the statistical uncertainty of any
PET result, and it is not paired with a PET central value; a larger ensemble
would sharpen the dispersion estimate and would still not supply coverage
evidence. These statements characterize the recoil-only estimator as method
development and do not supply a publication uncertainty for it or for a joint
muon--recoil estimator.

The historical recoil-PET campaign produced vertical, retraining, ML,
statistical and detector blocks on a common nominal.  That assembly is now
quarantined rather than quoted as an uncertainty.  For a nuisance endpoint
$u$, the physical shift and the induced retraining response form one joint
shift,
\begin{equation}
  \delta_u=x_u^{\rm varied+retrained}-x_{\rm CV}.
\end{equation}
Writing this as a frozen-map piece plus a retraining increment is an algebraic
decomposition, but adding the two separate covariances omits
$\operatorname{Cov}(\delta_{\rm frozen},\delta_{\rm retrain})$ and its
transpose.  Positive semidefiniteness of that sum does not restore the missing
cross term.

The historical detector block is additionally CV-support-limited: it propagates
shifted weights and observables through the nominal cloud membership instead of
regenerating the selected cloud for every lateral endpoint.  A replacement
construction would need end-to-end joint nuisance/retraining shifts and the
selection-complete FPS lateral samples, and would belong to a full-event PET
estimator with its own statistical and ML ensembles.  That is method-development
design and not a pending publication deliverable: no PET total covariance is
adopted for publication.  This section therefore quotes only the recoil-PET
central value and closure; no PET precision comparison is made and no PET
uncertainty magnitude is quoted.

\subsection{Truth-cloud projection within the recoil cross-check}
\label{sec:pet-projection}

The truth side of the recoil-only PET unfold is \emph{already} a point cloud: the
generator-level final-state hadron set
(\texttt{part\_gen}, the top-12-by-energy truth hadrons with their PDG codes,
built by \texttt{CVUniverse::GetTruthFSHadrons}) is carried per event
alongside the learned step-2 push weight. Any truth observable that is
a function of those 4-vectors can therefore be obtained after the
unfold by re-binning the push-weighted truth events --- no re-training and no
re-unfolding. We demonstrate this \emph{reweight-then-project} capability on
three observables computed post-hoc from the stored cloud: available energy
$\Eavail$, the hadronic-system invariant mass
$W_{\text{had}}=\sqrt{(\sum E)^2-|\sum\vec p|^2}$, and the proton/charged-pion
multiplicities. The exercise also quantifies the two
limitations of the stored cloud.

\paragraph{Truncation is essentially lossless.}
Because $\Eavail$ recomputed from the cloud uses the same formula and
inputs as the value stored at dump time, the cloud-vs-stored residual isolates
the cost of keeping only the 12 highest-energy hadrons. Over the full
\num{32.8}~million truth events (every event now carries a cloud, below) the
residual has median
\SI{0}{MeV} and mean \SI{\pcEavailMean}{MeV}, with \SI{\pcEavailWithin}{\percent}
of events within \SI{10}{MeV} (Fig.~\ref{fig:pcvalid}). The entire deficit
lives in the \SI{\pcSaturated}{\percent} of events that saturate the 12-particle
cap (median deficit \SI{\pcSatDeficit}{MeV} there, exactly zero below it), and
because the ranking is by energy the dropped constituents are the lowest-energy
stored final-state hadrons --- a small, one-sided, bounded effect. The 12-particle cap therefore
does not limit $\Eavail$-scale observables and does not need to be raised.

\paragraph{Coverage: closing the acceptance gap.}
Originally the cloud was filled only in the reco-signal event loop, so the
truth-only ``miss'' events appended for the unfold (\SI{\pcMiss}{\percent} of
the truth sample) carried no cloud --- only \SI{\pcHasCloudBefore}{\percent} of
truth events had one, and cloud projections were correspondingly
acceptance-limited, landing exactly on the stored-$\Eavail$ projection
restricted to the has-cloud subset but sitting $1.5$--$2\times$ below
the full unfold per bin. This is a truth-side instance of the same
dangling-miss-row dump issue tracked for the systematic bank. We fixed it at
source by carrying the truth final-state hadrons through the
truth-denominator loop and filling \texttt{part\_gen} for the appended miss
rows (the reco cloud \texttt{part\_reco} correctly stays empty for a miss),
then re-dumping, rebuilding the inputs, and re-training. Coverage is now
complete: \SI{\pcHasCloud}{\percent} of truth events carry a cloud (only
\pcEmptyCloud{} residual empties out of \num{32.8}~million), and the
cloud-recomputed $d\sigma/d\Eavail$ now reproduces the full unfold to
$\lesssim\SI{\pcCloudFull}{\percent}$ in every $\Eavail$ bin
(Fig.~\ref{fig:pcxsec}, center, where the cloud, stored, and full-unfold
curves coincide). The projection is no longer acceptance-limited.

\paragraph{Exploratory projection yields new hadronic diagnostics.}
$W_{\text{had}}$ is not a lossy copy of the leptonic ``experimenter's''
$W$ stored at dump time --- the latter is reconstructed from muon kinematics
assuming a single struck nucleon at rest, while $W_{\text{had}}$ is the true
invariant mass of the hadron system. The two differ by a median
\SI{\pcWoffset}{GeV}, and the 2D correlation
(Fig.~\ref{fig:pcxsec}, left) makes the physics explicit: a sharp locus at
$W_{\text{had}}\approx\SI{0.94}{GeV}$ (single-baryon final states, whose
hadronic mass is just the nucleon rest mass irrespective of leptonic $W$) and
diagonal bands stepped above it by $\sim$one nucleon mass per extra final-state
baryon (multi-nucleon knockout, 2p2h), Fermi-smeared. Projecting the cloud
thus provides a physically interpretable hadron-cloud diagnostic from the
existing recoil-only unfold; it is not promoted here as a separate full-event
measurement.

\begin{figure}[H]
  \centering
  \includegraphics[width=0.90\textwidth]{pet_cloud_projection_validation}
  \caption{Recoil-cloud truncation validation: cloud-recomputed $\Eavail$ minus the value
  stored at dump time, over the full \num{32.8}~million truth events.
  The residual is zero to \SI{\pcEavailWithin}{\percent} within \SI{10}{MeV};
  the small negative tail is confined to the \SI{\pcSaturated}{\percent} of
  events that reach the 12-hadron cap. The right panel shows ordinary event
  counts versus the pre-truncation multiplicity from the variable-length ROOT
  branch. Multiplicities 1--12 are shown exactly; the sparse tail is grouped
  into the labeled 13--14, 15--17, 18--24, and $\geq25$ categories. The dashed
  boundary separates fully stored events from categories for which only the
  12 highest-energy hadrons are retained. The cap-reaching events have a median
  $\Eavail$ deficit of only
  \SI{\pcSatDeficit}{MeV}. Keeping the 12 highest-energy hadrons is therefore
  essentially lossless for $\Eavail$-scale observables.}
  \label{fig:pcvalid}
\end{figure}

\begin{figure}[H]
  \centering
  \includegraphics[width=0.98\textwidth]{pet_cloud_projection_xsec}
  \caption{Exploratory recoil-only reweight-then-project demonstration.
  \emph{Left:} hadronic-system
  invariant mass $W_{\text{had}}$ versus the leptonic stored $W$ --- they
  differ by definition (median \SI{\pcWoffset}{GeV}), with the QE locus at
  \SI{0.94}{GeV} and multi-baryon bands above $y=x$. \emph{Center:}
  projected $d\sigma/d\Eavail$ from the push-weighted cloud now reproduces the
  full unfold to $\lesssim\SI{\pcCloudFull}{\percent}$ per bin after the
  miss-row coverage fix --- the cloud, stored-$\Eavail$, and full-unfold
  curves coincide; the GBDT 5D reference is shown for comparison. \emph{Right:}
  projected $d\sigma/d W_{\text{had}}$, whose shape differs from the
  leptonic-$W$ reference as expected.}
  \label{fig:pcxsec}
\end{figure}

\paragraph{The projection and extrapolation limits.}
\label{sec:pet-projection-next}
Reweighting can only redistribute weight among truth configurations the
generator already produced; it cannot populate phase space the prior never
sampled --- the same extrapolation limit flagged for the extended-fiducial
study (\S\ref{sec:fps}). Within that limit, hadron-cloud diagnostics can be
projected over the stored truth spectrum. They do not
turn the recoil-derived weight into a full-event weight or remove the prior
conditional for omitted muon information.
Two extensions lie beyond re-projection from the stored cloud (they would
require regenerating the event dump):
adding the muon through separately declared reco/data and truth event-feature
interfaces (it is excluded from the recoil cloud but retained as a stored
scalar); and raising the 12-particle cap for the rare
very-high-multiplicity tail (unnecessary at $\Eavail$ scale per the validation
above). Full details and reproduction are in
\texttt{nd-unfolding/pet/POINTCLOUD\_PROJECTION.md}.

\subsection{Scope of the recoil-PET cross-check}
The cross-check establishes that the point-cloud infrastructure works on
real MINERvA recoil data: paired reco/truth clouds, GPU training, full-sample
reweighting, absolute extraction, coherent estimator bootstraps and targeted
response studies all operate end to end. It also provides a useful comparison
against the scalar result and supports exploratory post-hoc hadron-cloud
projections. It does \emph{not} establish simultaneous muon-plus-hadron
unfolding, because the muon was not a classifier input.

A full-event extension requires separate observable
\texttt{event\_reco}/\texttt{event\_data} and truth-level
\texttt{event\_truth} interfaces containing a distinguished muon, explicit
masks/types and validated neighborhood geometry. An omitted-variable stress
closure that changes muon kinematics within narrow recoil strata would test
whether the joint estimator recovers information that the recoil-only
estimator necessarily misses. Any such estimator requires its own nominal fit
and uncertainty budget.
